\SourceLeaf{379}%
greater care was taken in every respect. The operations linked Mont Cenis with Mont Colombier by gunpowder signals on Mont Tabor, and Mont Colombier with the French triangulation point Puy d’Usson by gunpowder signals on the mountain Pierre sur autre. At the same time, signals were also given on Mont Colombier itself to link it with the Geneva observatory. The determination of time on Mont Cenis and Mont-Colombier was based on observations with transit telescopes by \textsc{Lenoir}\footnote{Judging by some of the circumstances, the same transit telescope seems to have been used on Mont Cenis in 1821 and 1822, although a different maker is named here.} and \textsc{Grindel}. At the French station it was based on absolute altitudes of stars measured with a repeating circle (see \textit{Conn. des tems} 1829 and our notice of 10 January 1828). On Mont Cenis they had to go three hours' journey from the hospice each time in order to see the signals. Finally, we also find here the determination of the difference in longitude between the observatories in Turin and Milan by gunpowder signals on S. Bernardo di Fenera on three separate occasions \(1823\mathpunct{.}\mathpunct{.}\mathpunct{.}1824\), with all the circumstances as favourable as they can be in operations of this kind.

The third section contains the determinations of latitude of the observatories on Mont Cenis, on Mont Colombier and in Turin, the first two with repeating circles by \textsc{Troughton} and \textsc{Reichenbach}, the last with the \textsc{Reichenbach} meridian circle. The latter observations do not show quite the degree of agreement to which one is usually accustomed with these instruments. The authors themselves pointed this out. and leave open whether such anomalies are real or depend on some error in handling the instrument that can be cleared up in the future. The reviewer recognizes that, given the many precautions this instrument requires, no one who was away from the site could offer even a plausible conjecture about this. But the way the authors discuss those anomalies prompts him to cite an example from his own experience of how seemingly trifling circumstances can sometimes impair observations. Over a number of years, the excellent agreement of observations with a meridian circle identical to the one at Turin was disturbed only once,
